% TOP_PROBLEM: {"id": 1101309, "title": "Questions in Geometric Group Theory — Q 13.9", "category_id": 17, "category": {"id": 17, "name": "group_theory", "display_name": "Group Theory", "description": "Problems about groups, group actions, representations, and related algebraic structures.", "slug": "group-theory", "order_index": 17, "created_at": "2026-07-31T15:26:25.671Z"}, "status": "open", "problem_number": "AMR-010-1309", "set_id": 13}
% FIRST_POSTED: 2026-10-02
% ENTRY_KIND: statement_only
% UnsolvedMath ID 1101309; source catalogue code AMR-010-1309
% !TeX program = lualatex
% Definition and sources only; this file is not an LLM solution attempt.
\documentclass[11pt,a4paper]{article}
\usepackage[margin=25mm]{geometry}
\usepackage{fontspec}
\setmainfont{lmroman10-regular.otf}[BoldFont=lmroman10-bold.otf,
 ItalicFont=lmroman10-italic.otf,BoldItalicFont=lmroman10-bolditalic.otf]
\usepackage{amsmath,amssymb,mathtools}
\usepackage{unicode-math}
\setmathfont{latinmodern-math.otf}
\AtBeginDocument{\let\setminus\smallsetminus}
\usepackage{xurl,hyperref}
\hypersetup{colorlinks=true,urlcolor=blue}
\setcounter{secnumdepth}{0}
\setlength{\parindent}{0pt}
\setlength{\parskip}{5pt}
\setlength{\emergencystretch}{3em}
\begin{document}
\section*{Questions in Geometric Group Theory — Q 13.9}\label{1101309}
{\small UnsolvedMath ID 1101309 · AMR-010-1309 · Group Theory · AMR Open Problem Lists}
\subsection{Definitions and mathematical statement}
For a finite presentation $P=\langle x_1,\ldots,x_m\mid r_1,\ldots,r_m\rangle$, write $G(P)$ for the presented group. Such a presentation is balanced because it uses equally many generators and relators. It presents a perfect group when $G(P)/[G(P),G(P)]=1$, equivalently when the exponent-sum matrix of its relators has determinant $1$ or $-1$.

Is there an effectively produced sequence $(P_j)_{j\geq1}$ of finite balanced presentations of perfect groups such that
\[
 T=\{j\geq1:G(P_j)=1\}
\]
is not a decidable subset of the positive integers? Effective production means an algorithm receives $j$ and returns the finite strings specifying $P_j$. Nondecidability means no algorithm halts on every $j$ and correctly tests membership in $T$.

The number of generators and the lengths of the relators may vary with $j$; the source does not impose a uniform bound on either. The perfectness promise applies to every member of the sequence, including the nontrivial groups, and prevents abelianization from deciding triviality in the intended examples.

This question concerns recognition of the trivial group across a family of presentations. It is different from the word problem inside a single group. A presentation being trivial can be verified by enumerating consequences until each generator is proved trivial, but the difficulty is obtaining an algorithm that also terminates on all the nontrivial members.
\subsection{Short English statement}
Can triviality be undecidable even within an effective sequence of balanced presentations of perfect groups?
\subsection{Sources}
\begin{itemize}
\item[S1] ULAM.AI and the UnsolvedMath Contributors, supplied problems.json, record 1101309 (AMR-010-1309), AMR Open Problem Lists. Catalogue identity and source transcription; CC BY 4.0.
\url{https://huggingface.co/datasets/ulamai/UnsolvedMath}.
\item[S2] Bestvina - Questions in Geometric Group Theory (pdf) (2004), Question 13.9 (PDF page 25). Original source indexed by AMR-010-1309.
\url{https://www.math.utah.edu/~bestvina/eprints/questions-updated.pdf}.
\end{itemize}
\end{document}
